\documentclass{article}
\usepackage{standalone}
\usepackage{hyperref}
\usepackage{datetime}
\usepackage{graphicx}
\usepackage{listings}
\usepackage{color}
\usepackage{todonotes}
\usepackage{lmodern}
\usepackage[T1]{fontenc}
\usepackage{scrextend}
\addtokomafont{labelinglabel}{\sffamily}
\usepackage{blindtext} 
\usepackage{marginnote}
\usepackage{hyperref}
\usepackage{smartdiagram}
\hypersetup{
    colorlinks=true,
    linktoc=all,
    linkcolor=blue,
    urlcolor=blue,
}
\author{Peter Mezei}
\title{Gardenzilla \\ development notes}
\pagestyle{myheadings}
\markright{Gardenzilla\hfill Development Notes \hfill}
\begin{document}
\maketitle
\thispagestyle{empty}
\newpage
\thispagestyle{empty}
\begin{figure}
  \centering
  \includegraphics[width=3cm]{logo.png}
\end{figure}
\begin{center}
  {\textit  Gardenzilla development notes, questions and future development directions. The main goal of this document is to document the development steps and collect all the ideas for future evaluation.

    \vspace{1cm} All right reserved \the\year{}

    All the logos, images and all the other\\ used visual materials are copyrighted.
  }
\end{center}
\vspace{11cm}
\begin{center}
  {\textbf {\huge </>}}
\end{center}
\newpage
\thispagestyle{empty}
\tableofcontents
\newpage
\section{Introduction}

\subsection{License}

\marginnote{License \\GNU GPL\\ version 2}

We mainly use GNU/GPL version 2. This means this is a free software,
you can freely use it, and modify it. IMPORTANT! As GNU/GPL version 2
is a COPYLEFT license, YOU MUST USE THE SAME GNU/GPL2 LICENSE in your
own verion. If you modify any of the project, YOU MUST MAKE YOUR
UPDATED VERSION AVAILABLE ONLINE. As this project is licensed under
GNU/GPLv2, you most apply the same license and use it - or publish it
- under the same GNU/GPLv2 licesen. If you use this project as a smaller
part of a larger `product`
\footnote{Software, Service, or any kind of
  business outcome where you use this product}
YOU MUST APPLY THE SAME GNU/GPLv2 LICENSE FOR YOUR LARGER PRODUCT AS WELL.

\subsection{Technology}

We use Rust as the main programming language. It is extremely powerfull,
effective, blazangly fast and safe by design. No race condition, and
almost impossible to create any bug that compiles. No garbage collection,
no runtime dependencies.

Other core technolgoies inside the project:

\begin{labeling}{alligator}
  \item [Gnu/Linux] \hfill \\ main development and hosting environment
  \item [Rust] Main programming language
  \item [GnuMakefile] \hfill \\ Using for main build and cleanup staffs
  \item [LaTex] for documentaion
  \item [Nginx] for webserver
  \item [Docker] for container management
  \item [TypeScript] \hfill \\ for typesafe javascript editing
  \item [HTML] for web documents
  \item [CSS] for web styles
  \item [SASS] for better css creation
  \item [JavaScript] \hfill \\ for client development
  \item [Angular] as the main client library
  \item [JSON] as the main API communication format
\end{labeling}

\subsection{3rd party libraries}

There are plenty awesome 3rd party libreries that we heavily use.
We would like to thank you for all the developers from the open-source
community who and whose projects helped us a lot.
In the above you can see the major 3rd party libs we use:

\subsubsection{3rd party Rust libraries}

\begin{labeling}{alligator}
  \item [storaget] data object layer as DAL. In the core project we do
  not use database, instead we use rust data objects inside
  storaget::Pack<T> and storaget::VecPack<T>. Pack manages all the
  data sync between memory and disk.\footnote{Storaget is our inhouse sub-project. See: \url{https://github.com/mezeipetister/storaget}}
  \item [chrono] date library. Using to manage date and datetimes. In
  the core project we use Date<Utc> and DateTime<Utc>. Whenever possible
  we avoid using NaiveDate(s).
  \item [rand] for random number generation.
  \item [bcrypt] for HTTP auth TOKEN encryption.
  \item [lettre] for SMTP email sending.
  \item [nanoid] for ID generation.
  \item [rocket] web framework for REST API.
  \item [rocket\_cors] \hfill \\ cors management for REST API
  \item [serde] object serialization, deserialization
  \item [jwt] for JSON web token management
\end{labeling}

\subsubsection{3rd party NodeJS libraries}

\begin{labeling}{alligator}
  \item [Angular] core web framework
  \item [chart.js] displaying charts
  \item [ng2-chartjs] \hfill \\ chart.js -> Agnular integration
  \item [markdown.it] \hfill \\ markdown management
  \item [markdown] \hfill \\ markdown management
\end{labeling}

\newpage

% section MODULES
% ***************
\section{Modules}

\subsection{Customer}
\label{customer}

We manage customers as independent data objects, and manage a customer
object through its implemented methods, or via public functions inside
core::customer module.

\begin{figure}[!htb]
  \caption{Customer object}
  \begin{verbatim}
    pub struct Customer {
      id: String,
      related_users: Vec<String>,
      name: String,
      tax_number: String,
      address: InvoiceAddress,
      phone: String,
      email: String,
      date_created: DateTime<Utc>,
      created_by: String,
    }
  \end{verbatim}
\end{figure}

Where:

\begin{description}
  \item[id] \hfill
        \\ Customer unique identifier \todo {Describe and set the way we manage IDs.}
  \item[related\_users] \hfill
        \\ vector of related user ID. User and Customer objects are independent
        of each other, as we can have customer without a user, and we can have
        - at least in theory - a user without a customer. Customer can
        purchase, User can login and interact withing the client. To manage
        Customer <-> User relations, we store them as a related ID list in
        each other. Here in a Customer object thats why we have this
        related\_users vector.

  \item[name] \hfill
        \\ name as string of the customer

  \item[tax\_number] \hfill
        \\ tax number as string
  \item[address] \hfill
        \\ address as InvoiceAddress object \todo{display InoviceAddress}
  \item[phone] \hfill
        \\ phone number as string
  \item[email] \hfill
        \\ email as string
  \item[date\_created] \hfill
        \\ date as chrono::DateTime<Utc>
  \item[created\_by] \hfill
        \\ userid as string, who created this object \todo{describe datecreated
          and created by somewhere as we use them everywhere}
\end{description}

\subsubsection{Methods}

\begin{description}
  \item[new(\dots) -> Self] \hfill
        \\ creates a new instance of Customer. ID cannot be changed after
        creation
  \item[get\_id() -> \&str] \hfill
        \\ returns customer id as \&str
  \item[get\_name() -> \&str] \hfill
        \\ returns customer name as \&str
  \item[set\_name(name: String)] \hfill
        \\ set customer name to the given new one
  \item[has\_user() -> bool] \hfill
        \\ returns bool if there is at least one related user
  \item[get\_users() -> \&Vec<String>] \hfill
        \\ returns related user vector reference
  \item[remove\_user(userid: \&str)] \todo{should return bool, or result} \hfill
        \\ remove a user by its userid
  \item[get\_tax\_number() -> \&str] \hfill
        \\ returns tax number as \&str
  \item[set\_tax\_number(tax\_number: String)] \hfill
        \\ set tax number
\end{description}

\newpage

\subsection{User}

At this first version, we have users only as admin users. As the online store
is not available, we do not need role management. All the registered users are
allowed to login.

Customers are also the part of the system, as the core object of the purchase process.
Currently all the purchase and invoice staff are related to customers. See at \ref{customer}.

\begin{figure}[!htb]
  \caption{User object}
  \begin{verbatim}
    pub struct User {
      id: String,                       // String, u32 or UserId?
      related_customers: Vec<String>,   // String, u32 or CustomerId?
      email: String,
      password_hash: String,            // Should we use PasswordHash instead?
      date_created: DateTime<Utc>,
      created_by: String,
    }
  \end{verbatim}
\end{figure}

Like:

\begin{verbatim}
  struct UserId(String);
  
  impl UserId {
    fn from(user_store: VecPack<User>, userid: &str) -> Self;
  }
\end{verbatim}

\begin{tabular}{|p{10cm}}
  \textbf{Fontos}
  Meg kell oldani, hogy a user-ben tárolt customerek és a customerben tárolt userek
  szinkronban legyenek egymással. Erre megoldás lehet, ha csak úgy engedünk hozzá adni,
  vagy elvenni egyikből-másikból, hogy paraméternek beadjuk a párt. Így az algoritmus felelhet
  azért, hogy a két entitás szinkronban legyen.

  Tehát
\end{tabular}

\begin{verbatim}
  let mut user = User::new();
  let mut customer = Customer::new();

  user.add_customer(&mut customer);
    |
    *---->  this should add the customer to the user as related user
            and also add user as related user to the given customer.

  user.remove_customer(&mut customer);
    |
    *---->  this should remove the related customer and also clean the current user
            from the given customer as well.
\end{verbatim}

\subsection{Product}

\begin{figure}[!htb]
  \caption{Product object}
  \begin{verbatim}
    pub struct Product {
      id: String,
      name: String,
      description: String,
      manufacturer: String,
      date_created: DateTime<Utc>,
      created_by: String,
    }
  \end{verbatim}
\end{figure}

\begin{quote}
  {\large \textbf{Q}uestions:}
  \begin{itemize}
    \item How to manage custom product data? - such as bee hazard level,
          or product hazard category? E.g.: in sales we need to know if a
          product requires any specific license or permission. But there can
          be many variety per product. E.g.: `gázmesteri, I-II-III forgalmi
          kategória növszereknél`. We also have different process per specific
          requirements: we need the permission license number for category
          II-I, but we need prescription for category I. We also need
          prescription for procurement of category I product.
    \item Best before date management? We have many products with 1-2-3
          years of expiration date, and we have many seasonal products, so
          if the next season start is later then the best before date, the real
          best before date is the end of the current season.
    \item Product related data? Such as active ingredient of a pesticide,
          or permitted dose per crops per product.
  \end{itemize}
\end{quote}

\subsection{Manufacturer}

\begin{figure}[!htb]
  \caption{Manufacturer object}
  \begin{verbatim}
    pub struct Manufacturer {
      id: String,
      name: String,
      date_created: DateTime<Utc>,
      created_by: String,
    }
  \end{verbatim}
\end{figure}

\newpage
\subsection{Partner}

Partner can be a source of our product, like a Wholesale or Warehouse.

\begin{quote}
  {\large \textbf{Q:}} Should
  \todo{Decide how to manage partner and customers}
  we manage our clients as partners as well?
  Or manage both of them separately Partners and Customers?
  \begin{itemize}
    \item A partner even can purchase from us? Yes, why not \dots
    \item A customer can be a partner? Theoretically yes.
    \item Do we wanna see our customers in our partner list? Probably not.
    \item Do we wanna see our partners in our customer list? Probably we don't mind.
  \end{itemize}
\end{quote}

\begin{figure}[!htb]
  \caption{Partner object}
  \begin{verbatim}
    pub struct Partner {
      id: String,
      name: String,
      description: String,
      manufacturer: String,
      date_created: DateTime<Utc>,
      created_by: String,
    }
  \end{verbatim}
\end{figure}

\newpage

\subsection{Cash}

\begin{verbatim}
    pub struct CashRegister {
      balance: i32,
      transaction: Vec<Transaction>
    }
  \end{verbatim}

\begin{verbatim}
    pub struct Transaction {
      id: String,
      name: String,
      amount: i32,
      kind: TransactionKind
      created_by: String,
      date_created: DateTime<Utc>,
    }
  \end{verbatim}

\begin{verbatim}
    pub enum TransactionKind {
      MoneyIn {
        who: String,
        amount: i32,
        comment: String,
      },
      MoneyOut {
        who: String,
        amount: i32,
        comment: String,
      },
      Purchase {
        purchase_id: String,
      },
      Change {
        comment: String,
      },
    }
  \end{verbatim}

\subsection{Reservation}

\subsection{Stock}

\begin{verbatim}
  pub struct Stock {
    id: String,
    products: Vec<StockItem>
  }

  pub struct StocSku {
    sku: String,
    items: Vec<StockItem>
  }

  pub struct StockItem {
    upid: String,
  }
\end{verbatim}

We\todo{Unique Product Layer (UPL?)} should use a spacial product layer, to hold all the unique product
as a single product holder, with alll of its data and history; - and store
this product holder inside the stock product vector.

\begin{quote}
  Questions:
  \begin{itemize}
    \item It should be really fast to check a UPL by ID.
    \item It should be easy to locate where the UPLs are located.
          And their movements.
    \item We should hold the minimum required data,
          and manage memory smartly - as we must estimate a large number of
          UPL objects.
  \end{itemize}
\end{quote}

\subsection{Procurement}

\begin{lstlisting}[caption='Procurement Object']
  pub struct Procurement {
    
  }
\end{lstlisting}

\begin{quote}
  Questions:
  \begin{itemize}
    \item It should be really fast to check a UPL by ID.
    \item It should be easy to locate where the UPLs are located.
          And their movements.
    \item We should hold the minimum required data,
          and manage memory smartly - as we must estimate a large number of
          UPL objects.
  \end{itemize}
\end{quote}

\subsection{Source}

\begin{lstlisting}[caption='Source Object']
  type Sku = String;
  
  pub struct Source {
    id: String,
    name: String,
    items: HashMap<Sku, ItemHistory>
  }

  pub struct ItemHistory {
    price: i32,
    created_by: String
    date_created: DateTime<Utc>,
  }
\end{lstlisting}

\subsection{Purchase}

\subsection{Cart}

\subsection{Issue}

\subsection{Invoice}


\subsection{Document?}

\subsection{Customer Card}

\begin{lstlisting}
  pub struct CustomerCard {
    id: String,
    owner: CustomerId,
    discount_pct: u32,
    date_created: DateTime<Utc>,
    created_by: UserId
  }
\end{lstlisting}

Kérdések

\begin{enumerate}
  \item ez minden customer-nek jár?
  \item egy customernek csak egy lehet?
  \item ez a customer profil része, vagy külön modul?
  \item lehet ezt a kártyát cserélni?
  \item online vásárlásnál automatikusan megy rá?
  \item offline vásárlásnál kell a fizikai kártya?
  \item history legyen benne? Vagy a CustomerId alapján kereshető?
  \item Mi a célja? Elköteleződés, illetve az offline beazonosítás.
\end{enumerate}

\documentclass{article}
\usepackage{listings}
\begin{document}
\subsection{Gift Card}

\begin{lstlisting}[caption='GiftCard Object']
  pub struct GiftCard {
    id: String,
    owner: CustomerId,
    balance: u32,
    history: Vec<Transaction>
  }

  pub struct Transaction {
    id: i32,
    amount: i32,
    reference: Reference,
    date_created: DateTime<Utc>,
    created_by: UserId,
  }

  pub enum Reference {
    Purchase(PurchaseId),
    MoneyIn(CashId),
    MoneyBack(CashId)
  }
\end{lstlisting}
\end{document}

Hello \makebox[5cm]{\dotfill} you are welcomed. Me as \makebox[5cm]{\dotfill} aggrees that \dots

% A circular diagram of a TeX workflow
% Author: Stefan Kottwitz
% https://www.packtpub.com/hardware-and-creative/latex-cookbook
\documentclass{standalone} 
\usepackage{verbatim}
\begin{comment}
:Title: A circular diagram of a TeX workflow
:Tags: Diagrams;Flowcharts;Smartdiagram;Cookbook
:Author: Stefan Kottwitz
:Slug: smart-circle

A diagram showing a TeX workflow using the
clever smartdiagram package.
\end{comment}
\usepackage{smartdiagram}
\begin{document}
\smartdiagram[circular diagram:clockwise]{Landing page,
  Attention management, Target page, Cart/Checkout, Welcome}
\end{document}

\newpage
\todototoc
\listoftodos

\newpage
\thispagestyle{empty}
\vspace*{5cm}
\begin{center}
  {\textit Gardenzilla\\ Development Notes \\ \vspace{5mm}
    All right reserved \\ Peter Mezei \\\the\year{}}
\end{center}
\end{document}
